\documentclass[12pt,a4paper]{article}

% -------------------- Packages --------------------
\usepackage{graphicx}
\usepackage{float}
\usepackage{amsmath,amssymb,mathtools}
\usepackage{geometry}
\usepackage{xcolor}
\usepackage{booktabs}
\usepackage{caption}
\usepackage{array}
\usepackage{tabularx}
\usepackage{enumitem}
\usepackage{titlesec}
\usepackage{fancyhdr}
\usepackage{hyperref}
\usepackage[most]{tcolorbox}

% XePersian should usually be loaded near the end.
\usepackage{xepersian}

% -------------------- Page Layout --------------------
\geometry{
	top=2.5cm,
	bottom=2.5cm,
	left=2.5cm,
	right=2.5cm
}

% -------------------- Fonts --------------------
\settextfont{XB Niloofar}
\setlatintextfont{Times New Roman}

% -------------------- Spacing --------------------
\linespread{1.2}
\setlength{\parindent}{1.2em}
\setlength{\parskip}{0.35em}

% -------------------- Colors --------------------
\definecolor{mainblue}{RGB}{35, 70, 120}
\definecolor{softgray}{RGB}{245, 245, 245}

% -------------------- Section Styling --------------------
\titleformat{\section}
{\Large\bfseries\color{mainblue}}
{\thesection}
{0.8em}
{}

\titleformat{\subsection}
{\large\bfseries\color{mainblue}}
{\thesubsection}
{0.7em}
{}

\titlespacing*{\section}{0pt}{1.4em}{0.7em}
\titlespacing*{\subsection}{0pt}{1.1em}{0.5em}

% -------------------- Lists --------------------
\setlist[itemize]{
	topsep=0.3em,
	itemsep=0.25em,
	parsep=0pt,
	rightmargin=1.5em
}

\setlist[enumerate]{
	topsep=0.3em,
	itemsep=0.3em,
	parsep=0pt,
	rightmargin=1.5em
}

% -------------------- Captions --------------------
\captionsetup{
	font=small,
	labelfont=bf,
	justification=centering,
	skip=8pt
}

% -------------------- Tables --------------------
\renewcommand{\arraystretch}{1.25}
\newcolumntype{Y}{>{\centering\arraybackslash}X}

% -------------------- Header / Footer --------------------
\pagestyle{fancy}
\fancyhf{}
\fancyhead[R]{گزارش پروژه}
\fancyhead[L]{مدل‌سازی آماری و مونت‌کارلو}
\fancyfoot[C]{\thepage}
\renewcommand{\headrulewidth}{0.4pt}

% -------------------- Hyperlinks --------------------
\hypersetup{
	colorlinks=true,
	linkcolor=mainblue,
	citecolor=mainblue,
	urlcolor=mainblue
}



\begin{document}
	\begin{titlepage}
		\thispagestyle{empty}
		\centering
		
		\vspace*{1cm}
		
		{\Large\bfseries گزارش پروژه\par}
		
		\vspace{0.45cm}
		
		{\color{mainblue}\rule{0.78\textwidth}{1.2pt}}
		
		\vspace{1.1cm}
		
		{\Huge\bfseries\color{mainblue}
			مدل‌سازی آماری و تحلیل رویدادهای نادر
			\par
		}
		
		\vspace{0.35cm}
		
		{\LARGE\bfseries
			با روش‌های پیشرفته مونت‌کارلو
			\par
		}
		
		\vspace{4.1cm}
		
		\begin{center}
			\begin{tabular}{rl}
				\textbf{نام و نام خانوادگی:} & علی تقی‌زاده \\[0.3cm]
				\textbf{شماره دانشجویی:}     & ۴۰۳۱۳۰۰۸
			\end{tabular}
		\end{center}
		
		\vspace{1cm}
		
		\begin{tcolorbox}[
			width=0.82\textwidth,
			colback=white,
			colframe=mainblue,
			arc=4mm,
			boxrule=0.8pt,
			left=8mm,
			right=8mm,
			top=5mm,
			bottom=5mm
			]
			\centering
			
			{\large\bfseries مخزن کد پروژه\par}
			
			\vspace{0.4cm}
			
			\lr{
				\href{https://github.com/alitkbbl/Water-Quality-Rare-Events}
				{github.com/alitkbbl/Water-Quality-Rare-Events}
			}
		\end{tcolorbox}
		
		\vfill
		
		{\large
			تحلیل کیفیت آب، برازش آماری، شبیه‌سازی مونت‌کارلو و نمونه‌گیری اهمیتی
			\par
		}
		
		\vspace{0.45cm}
		
		{\color{mainblue}\rule{0.55\textwidth}{0.8pt}}
		
		\vspace{0.4cm}
		
		{\large\today\par}
	\end{titlepage}
	
	
	\setcounter{page}{1}
	
	
	
	\section{مقدمه}


در تحلیل سیستم‌های حساس، مانند بررسی کیفیت آب، تنها دانستن وضعیت میانگین یا رفتار معمول سیستم کافی نیست. در بسیاری از کاربردهای واقعی، آنچه اهمیت بیشتری دارد، پیش‌بینی احتمال وقوع شرایط بحرانی است؛ برای نمونه، حالتی که غلظت یک ماده شیمیایی از حد مجاز فراتر می‌رود. چنین رخدادهایی ممکن است به‌ندرت اتفاق بیفتند، اما اثر آن‌ها می‌تواند بسیار مهم و حتی خطرناک باشد.

در این پروژه، تمرکز اصلی بر تحلیل غلظت کلرامین \lr{(Chloramines)} قرار دارد. روند کار از شناخت اولیه داده‌ها و استخراج یک مدل آماری مناسب آغاز شد. سپس برای شبیه‌سازی داده‌ها، روش پایه رد-پذیرش پیاده‌سازی شد. در ادامه، مسئله اصلی پروژه یعنی محاسبه احتمال یک رویداد نادر بررسی شد و نشان داده شد که روش مونت‌کارلو استاندارد \lr{(Standard Monte Carlo)} در چنین شرایطی کارایی مناسبی ندارد. در نهایت، با استفاده از روش نمونه‌گیری اهمیت \lr{(Importance Sampling)}، کاهش چشمگیری در هزینه محاسباتی و خطای تخمین به دست آمد.

\section{شناخت داده‌ها و استخراج مدل آماری پایه}

اولین گام در هر پروژه شبیه‌سازی، شناخت رفتار واقعی داده‌هاست. در این پروژه، کار با مجموعه داده‌ای شامل $3276$ مشاهده معتبر از غلظت کلرامین آغاز شد. تحلیل‌های اولیه نشان داد که میانگین غلظت برابر با $7.122$ و واریانس آن برابر با $2.505$ است.

برای شبیه‌سازی این متغیر در محیط کامپیوتر، لازم بود یک مدل ریاضی مناسب برای تابع چگالی احتمال، یا همان \lr{PDF}، انتخاب شود. از آنجا که غلظت یک ماده شیمیایی نمی‌تواند مقدار منفی داشته باشد، دو توزیع گاما و وایبول به عنوان گزینه‌های اصلی در نظر گرفته شدند. سپس با استفاده از روش برآورد بیشینه درست‌نمایی \lr{(MLE)}، پارامترهای بهینه هر توزیع محاسبه شد:

\begin{table}[H]
	\centering
	\caption{پارامترهای برآوردشده برای توزیع‌های کاندید}
	\label{tab:distribution-parameters}
	\begin{tabularx}{0.85\textwidth}{Y Y Y}
		\toprule
		\textbf{توزیع} & \textbf{پارامتر اول} & \textbf{پارامتر دوم} \\
		\midrule
		گاما & شکل: $k = 18.0283$ & مقیاس: $\theta = 0.3951$ \\
		وایبول & شکل: $c = 4.8842$ & مقیاس: $\lambda = 7.7437$ \\
		\bottomrule
	\end{tabularx}
\end{table}


\begin{itemize}
	\item توزیع گاما: پارامتر شکل $k = 18.0283$ و پارامتر مقیاس $\theta = 0.3951$
	\item توزیع وایبول: پارامتر شکل $c = 4.8842$ و پارامتر مقیاس $\lambda = 7.7437$
\end{itemize}



پرسش اصلی این بود که کدام‌یک از این دو مدل، رفتار واقعی داده‌ها را بهتر نمایش می‌دهد.

\begin{figure}[H]
	\centering
	\includegraphics[width=0.85\textwidth]{"../results/01_Chloramines_distribution_fit.png"}
	\caption{برازش توزیع‌ها روی هیستوگرام داده‌های کلرامین}
	\label{fig:distribution-fitting}
\end{figure}

\subsection{تحلیل نمودار و انتخاب مدل نهایی}

همان‌طور که در شکل \ref{fig:distribution-fitting} مشاهده می‌شود، توزیع وایبول، که با خط‌چین سبز نمایش داده شده است، در تخمین نقطه اوج داده‌ها دقت کافی ندارد و قله توزیع را کمتر از مقدار واقعی نشان می‌دهد. در مقابل، توزیع گاما، که با خط قرمز پیوسته مشخص شده است، هم در بخش مرکزی داده‌ها و هم در نواحی ابتدایی و انتهایی نمودار، انطباق بهتری با هیستوگرام تجربی دارد.

بنابراین، توزیع گاما به عنوان مدل آماری اصلی و تابع هدف پروژه انتخاب شد. از این مرحله به بعد، این تابع با نماد $f(x)$ نمایش داده می‌شود و مبنای شبیه‌سازی‌ها و تحلیل‌های بعدی خواهد بود.

بنابراین تابع هدف نهایی به صورت زیر است:
\begin{equation}
	f(x) = \frac{1}{\Gamma(k)\,\theta^{k}}\, x^{k-1} e^{-x/\theta},
	\qquad k = 18.0283,\; \theta = 0.3951.
\end{equation}


\section{پیاده‌سازی نمونه‌گیری با الگوریتم رد-پذیرش}

پس از انتخاب تابع هدف $f(x)$، مسئله بعدی این بود که چگونه می‌توان اعداد تصادفی‌ای تولید کرد که از همین توزیع گاما پیروی کنند. برای درک بهتر سازوکار نمونه‌گیری، به جای استفاده از توابع آماده، الگوریتم رد-پذیرش \lr{(Accept-Reject)} از ابتدا پیاده‌سازی شد.

ایده اصلی این روش آن است که از یک توزیع ساده‌تر به عنوان توزیع پیشنهادی استفاده شود. در این پروژه، توزیع یکنواخت روی بازه $[0,16]$ انتخاب شد. بنابراین تابع چگالی پیشنهادی برابر است با:

\[
g(x) = 0.0625
\]

چالش اصلی در این روش، یافتن ضریب $M$ است؛ به گونه‌ای که برای تمام مقادیر $x$ داشته باشیم:

\[
M \cdot g(x) \geq f(x)
\]

با محاسبه نقطه بیشینه توزیع گاما مشخص شد که مقدار بیشینه در $x = 6.7279$ رخ می‌دهد و مقدار تابع در این نقطه برابر با $0.2435$ است. در نتیجه، ضریب $M$ به صورت زیر به دست آمد:

\[
M = 3.8959
\]

\begin{figure}[H]
	\centering
	\includegraphics[width=0.85\textwidth]{"../results/02_accept_reject_results.png"}
	\caption{نتایج نمونه‌گیری با الگوریتم رد-پذیرش}
	\label{fig:accept-reject}
\end{figure}

\subsection{تحلیل عملکرد الگوریتم}

الگوریتم برای تولید $10000$ نمونه معتبر اجرا شد. در عمل، برنامه مجبور شد $39709$ نمونه کاندید تولید کند تا در نهایت تعداد نمونه‌های پذیرفته‌شده به مقدار مورد نظر برسد. نرخ پذیرش تجربی برابر با $0.2518$، یعنی حدود $25$ درصد، به دست آمد.

این نتیجه با مقدار تئوری نیز سازگار است. از نظر نظری، نرخ پذیرش مورد انتظار برابر است با:

\[
\frac{1}{M} = 0.2567
\]

نزدیکی مقدار تجربی $0.2518$ به مقدار تئوری $0.2567$ نشان می‌دهد که پیاده‌سازی الگوریتم به‌درستی انجام شده است. همچنین، همان‌طور که در شکل \ref{fig:accept-reject} دیده می‌شود، هیستوگرام نمونه‌های پذیرفته‌شده به‌خوبی با منحنی توزیع هدف منطبق است. این موضوع تأیید می‌کند که نمونه‌های تولیدشده واقعاً از توزیع گامای مورد نظر پیروی می‌کنند.

\section{شبیه‌سازی مونت‌کارلو استاندارد و چالش رویدادهای نادر}

هدف نهایی پروژه، محاسبه احتمال وقوع یک وضعیت بحرانی بود؛ یعنی احتمال اینکه غلظت کلرامین از آستانه مشخص $\tau = 12.0$ فراتر رود. مقدار واقعی و تحلیلی این احتمال برابر است با:

\[
\gamma = 0.00621
\]

این مقدار تقریباً معادل $0.62\%$ است. به بیان ساده‌تر، چنین رخدادی به طور میانگین در حدود $6$ بار از هر $1000$ نمونه اتفاق می‌افتد. بنابراین با یک رویداد نادر روبه‌رو هستیم.

در ابتدا، روش مونت‌کارلو استاندارد \lr{(SMC)} برای تخمین این احتمال به کار گرفته شد. نتایج اولیه نشان دادند که این روش در حجم نمونه‌های پایین، دقت مناسبی ندارد:

\begin{itemize}
	\item برای حجم نمونه $N = 1000$، مقدار احتمال برابر با $0.002$ تخمین زده شد. خطای نسبی یا \lr{CV} در این حالت برابر با $0.7064$ بود؛ یعنی حدود $70$ درصد خطا.
	\item برای رساندن خطا به کمتر از یک درصد، یعنی $CV < 0.01$، لازم بود حجم نمونه تا $N = 10000000$ افزایش یابد. در این حالت مقدار \lr{CV} به حدود $0.004$ رسید.
\end{itemize}

\begin{figure}[H]
	\centering
	\includegraphics[width=0.9\textwidth]{"../results/03_standard_mc_results.png"}
	\caption{همگرایی روش مونت‌کارلو استاندارد و رفتار خطای نسبی در مقیاس لگاریتمی}
	\label{fig:smc-convergence}
\end{figure}

\subsection{تحلیل محدودیت‌های روش مونت‌کارلو استاندارد}


شکل \ref{fig:smc-convergence} نشان می‌دهد که خطای روش مونت‌کارلو استاندارد با نرخ تئوری زیر کاهش می‌یابد:

\[
\frac{1}{\sqrt{N}}
\]

این نرخ کاهش خطا، به‌خصوص در مسئله رویدادهای نادر، بسیار کند است. دلیل آن نیز روشن است: در روش استاندارد، بیشتر نمونه‌ها در ناحیه معمولی توزیع تولید می‌شوند؛ یعنی حوالی میانگین که در این مسئله نزدیک به مقدار $7$ است. اما رویداد مورد نظر ما در ناحیه بحرانی، یعنی مقادیر بزرگ‌تر از $12$، رخ می‌دهد.

در نتیجه، بخش بزرگی از زمان محاسباتی صرف تولید نمونه‌هایی می‌شود که هیچ نقشی در تخمین احتمال رویداد نادر ندارند. این مسئله باعث می‌شود روش مونت‌کارلو استاندارد برای چنین کاربردهایی بسیار پرهزینه و ناکارآمد باشد.

از نظر آماری، چون $\mathbb{1}_{\{X_i > \tau\}}$ یک متغیر برنولی با احتمال موفقیت $\gamma$ است، واریانس برآوردگر مونت‌کارلو استاندارد به صورت زیر به دست می‌آید:

\begin{equation}
	\mathrm{Var}\!\left(\hat{\gamma}_{\mathrm{SMC}}\right)
	=
	\frac{\gamma(1-\gamma)}{N},
	\qquad
	\mathrm{CV}
	=
	\frac{\sqrt{\mathrm{Var}\!\left(\hat{\gamma}_{\mathrm{SMC}}\right)}}{\gamma}
	\approx
	\sqrt{\frac{1-\gamma}{\gamma N}}.
	\label{eq:smc-variance}
\end{equation}

برای یک رویداد نادر، مقدار $\gamma$ بسیار کوچک است. در این پروژه داریم $\gamma \approx 6.21 \times 10^{-3}$؛ بنابراین خطای نسبی تقریباً به عبارت زیر وابسته می‌شود:

\begin{equation}
	\mathrm{CV} \approx \frac{1}{\sqrt{\gamma N}}.
	\label{eq:smc-cv-rare-event}
\end{equation}

رابطه \eqref{eq:smc-cv-rare-event} نشان می‌دهد که با کوچک شدن $\gamma$، خطای نسبی به‌سرعت افزایش می‌یابد. به همین دلیل، برای $N = 1000$ مقدار \lr{CV} حدود $0.71$ است و تخمین $0.002$ در برابر مقدار واقعی $0.00621$ دقت مناسبی ندارد. تنها زمانی به دقت قابل قبول می‌رسیم که حجم نمونه تا حدود $10^6$ یا $10^7$ افزایش یابد.

این پدیده همان مشکل واریانس بالا در شبیه‌سازی رویدادهای نادر است؛ یعنی برای رسیدن به یک سطح خطای ثابت، تعداد نمونه‌ها باید تقریباً متناسب با $\frac{1}{\gamma}$ رشد کند.




\section{رویکرد هوشمند با نمونه‌گیری اهمیت}

برای عبور از محدودیت روش مونت‌کارلو استاندارد، از تکنیک نمونه‌گیری اهمیت \lr{(Importance Sampling)} استفاده شد. ایده اصلی این روش آن است که به جای نمونه‌گیری کورکورانه از توزیع اصلی، نمونه‌ها از توزیعی تولید شوند که تمرکز بیشتری بر ناحیه مهم مسئله دارد.

در این پروژه، ناحیه مهم همان ناحیه بحرانی است؛ یعنی مقادیر بزرگ‌تر از آستانه $\tau = 12.0$. بنابراین یک توزیع پیشنهادی جدید با نماد $q(x)$ تعریف شد که جرم احتمالی بیشتری را در این ناحیه قرار می‌دهد.

برای متمرکز کردن نمونه‌ها در ناحیه بحرانی $x > \tau$، توزیع پیشنهادی $q(x)$ به صورت یک توزیع گاما انتخاب شد؛ با این تفاوت که میانگین آن به سمت آستانه بحرانی جابه‌جا شده است. این انتخاب باعث می‌شود قله توزیع پیشنهادی نزدیک $\tau = 12$ قرار گیرد و در نتیجه، نمونه‌های بیشتری از ناحیه مؤثر در محاسبه احتمال رویداد نادر تولید شوند.

برآوردگر نمونه‌گیری اهمیتی \lr{(Importance Sampling)} و وزن اهمیت به صورت زیر تعریف می‌شوند:

\begin{equation}
	\hat{\gamma}_{\mathrm{IS}}
	=
	\frac{1}{N}
	\sum_{i=1}^{N}
	\mathbb{1}_{\{X_i > \tau\}}\, w(X_i),
	\qquad
	w(x)
	=
	\frac{f(x)}{q(x)}.
	\label{eq:is-estimator}
\end{equation}

در رابطه \eqref{eq:is-estimator}، تابع $f(x)$ چگالی هدف، یعنی همان توزیع گامای برازش‌شده روی داده‌های کلرامین، و تابع $q(x)$ چگالی پیشنهادی است. وزن $w(x)$ نیز نقش جبرانی دارد و باعث می‌شود با وجود نمونه‌گیری از $q(x)$، تخمین نهایی همچنان نسبت به توزیع اصلی $f(x)$ معتبر باقی بماند.


\begin{figure}[H]
	\centering
	\includegraphics[width=0.85\textwidth]{"../results/04_IS_Strategy.png"}
	\caption{استراتژی نمونه‌گیری اهمیت با تمرکز روی ناحیه بحرانی}
	\label{fig:importance-sampling}
\end{figure}

\subsection{استراتژی و نتایج}

همان‌طور که در شکل \ref{fig:importance-sampling} مشخص است، توزیع پیشنهادی $q(x)$ به گونه‌ای انتخاب شده که قله آن نزدیک به ناحیه وقوع رویداد نادر قرار گیرد. به این ترتیب، نمونه‌های بیشتری در حوالی مقادیر بزرگ‌تر از $12$ تولید می‌شوند. البته برای اینکه تخمین همچنان بدون اریبی باقی بماند، از وزن‌های جبرانی یا نسبت درست‌نمایی \lr{(Likelihood Ratio)} استفاده شد.


نتایج به‌دست‌آمده بسیار قابل توجه بودند. روش نمونه‌گیری اهمیت تنها با $N = 1000$ نمونه توانست احتمال را برابر با $0.00612$ تخمین بزند و خطای نسبی را به مقدار $0.0387$ برساند. این مقدار کمتر از $4$ درصد است و برای چنین حجم نمونه‌ای بسیار مناسب محسوب می‌شود.

\begin{table}[H]
	\centering
	\caption{مقایسه عددی عملکرد روش مونت‌کارلو استاندارد و نمونه‌گیری اهمیت}
	\label{tab:smc-is-comparison}
	\begin{tabularx}{0.95\textwidth}{Y Y Y Y}
		\toprule
		\textbf{روش} & \textbf{حجم نمونه} & \textbf{تخمین احتمال} & \textbf{خطای نسبی} \\
		\midrule
		\lr{SMC} & $1000$ & $0.00200$ & $0.7064$ \\
		\lr{SMC} & $100000$ & $0.00621$ & $0.0400$ \\
		\lr{IS} & $1000$ & $0.00612$ & $0.0387$ \\
		\lr{IS} & $100000$ & نزدیک به مقدار تحلیلی & $0.0040$ \\
		\bottomrule
	\end{tabularx}
\end{table}


جدول \ref{tab:smc-is-comparison} نشان می‌دهد که روش نمونه‌گیری اهمیت با تنها $1000$ نمونه، به دقتی تقریباً مشابه روش مونت‌کارلو استاندارد با $100000$ نمونه می‌رسد. بنابراین، در این مسئله کاهش هزینه محاسباتی تقریباً صدبرابری حاصل شده است.


این مقایسه نشان می‌دهد که روش نمونه‌گیری اهمیت توانسته است با تعداد نمونه‌ای $100$ برابر کمتر، به دقتی تقریباً مشابه و حتی کمی بهتر برسد. همچنین در حجم نمونه $100000$، خطای روش \lr{IS} به مقدار بسیار کوچک $0.0040$ رسید؛ در حالی که روش \lr{SMC} برای رسیدن به چنین سطحی از دقت، به حدود $10000000$ نمونه نیاز داشت.

\section{جمع‌بندی نهایی و مقایسه کارایی محاسباتی}

پس از پیاده‌سازی و تحلیل جداگانه دو روش مونت‌کارلو استاندارد و نمونه‌گیری اهمیت، مرحله نهایی مقایسه مستقیم کارایی آن‌هاست. در مسائل مهندسی و علوم محاسباتی، یک روش خوب فقط نباید به پاسخ درست برسد؛ بلکه باید این کار را با هزینه محاسباتی قابل قبول انجام دهد.

در این پروژه، هزینه محاسباتی با تعداد نمونه‌های مورد نیاز، یعنی $N$، سنجیده می‌شود. همچنین کارایی روش‌ها بر اساس سرعت کاهش خطای نسبی یا \lr{CV} ارزیابی شده است.

\begin{figure}[H]
	\centering
	\includegraphics[width=0.9\textwidth]{"../results/05_SMC_vs_IS_Efficiency_Comparison.png"}
	\caption{مقایسه کارایی روش مونت‌کارلو استاندارد و نمونه‌گیری اهمیت}
	\label{fig:efficiency-comparison}
\end{figure}

\subsection{تحلیل نهایی}

شکل \ref{fig:efficiency-comparison} خلاصه اصلی نتایج پروژه را به‌خوبی نشان می‌دهد. در این نمودار، خطای نسبی هر دو روش در برابر حجم نمونه، در مقیاس لگاریتمی-لگاریتمی رسم شده است. فاصله زیاد بین دو منحنی نشان می‌دهد که نمونه‌گیری اهمیت از نظر محاسباتی بسیار کاراتر از روش مونت‌کارلو استاندارد است.

چند نکته اصلی از این نمودار قابل برداشت است:

\begin{enumerate}
	\item \textbf{شروع متفاوت دو روش:}
	روش \lr{SMC} در حجم نمونه پایین با خطای بسیار بزرگی آغاز می‌شود. برای مثال، در $N = 1000$، خطای نسبی آن حدود $70$ درصد است. در مقابل، روش \lr{IS} با همین تعداد نمونه، خطایی کمتر از $4$ درصد دارد.

	\item \textbf{کاهش خطا با افزایش حجم نمونه:}
	در هر دو روش، با افزایش $N$، خطا کاهش می‌یابد. با این حال، منحنی مربوط به \lr{IS} تقریباً در تمام نقاط، چندین مرتبه پایین‌تر از منحنی \lr{SMC} قرار دارد.

	\item \textbf{معنای عملی اختلاف دو روش:}
	برای رسیدن به خطای نسبی حدود $4$ درصد، یعنی $CV = 0.04$، روش \lr{SMC} به $N = 100000$ نمونه نیاز دارد. در حالی که روش \lr{IS} همین دقت را با $N = 1000$ نمونه به دست می‌آورد.
\end{enumerate}

بنابراین، با یک تغییر هوشمندانه در استراتژی نمونه‌گیری و تمرکز دادن نمونه‌ها روی ناحیه بحرانی، توانستیم همان سطح دقت را با کاهش صدبرابری هزینه محاسباتی به دست آوریم.

\section{نتیجه‌گیری}

این پروژه نشان داد که در تحلیل مسائل واقعی، به‌ویژه در حوزه‌هایی مانند تحلیل ریسک، قابلیت اطمینان و پایش کیفیت، رویدادهای نادر نقش بسیار مهمی دارند. اگرچه روش مونت‌کارلو استاندارد از نظر مفهومی ساده و قابل فهم است، اما در مواجهه با رویدادهای نادر می‌تواند بسیار کند و پرهزینه باشد.

در مقابل، نمونه‌گیری اهمیت \lr{(Importance Sampling)} نشان داد که با انتخاب مناسب توزیع پیشنهادی و استفاده از وزن‌های جبرانی، می‌توان دقت تخمین را به شکل قابل توجهی افزایش داد. این روش صرفاً یک بهینه‌سازی کوچک نیست، بلکه در چنین مسئله‌هایی ابزاری ضروری برای تبدیل یک محاسبه سنگین و زمان‌بر به یک تحلیل سریع، دقیق و قابل اعتماد است.

\end{document}
